\documentclass[../../../include/open-logic-section]{subfiles}

\begin{document}

\olfileid{sth}{replacement}{strength}
\olsection{ప్రతిస్థాపన బలం}

ప్రతిస్థాపన బలం గురించి ఒక సరళమైన పరిశీలనతో మొదలెడదాం: $\Z$ను దాటి వెళ్లకపోతే, $\omega + \omega$కు సమానమైన లేదా అంతకంటే పెద్ద వాన్ నాయ్‌మన్ క్రమసంఖ్య ఏదైనా ఉందని నిరూపించలేం.

ఎందుకో ఒక రూపురేఖ చూద్దాం. $\ZF$లో పనిచేస్తూ $V_{\omega+\omega}$ సమితిని పరిగణించండి. ఈ సమితి $\Z$కు ఒక \emph{నమూనా} యొక్క నిర్వచన సమితిగా పనిచేస్తుంది. ఇది చూడడానికి, సూత్రాన్ని \emph{పరిధికి పరిమితం చేయడం} కోసం ఒక సంకేతీకరణను పరిచయం చేద్దాం:
\begin{defn}\ollabel{formularelativization}
	ఏ సమితి $M$, ఏ సూత్రం $\phi$కైనా, $\phi$లోని అన్ని పరిమాణకారకాలను $M$కు పరిమితం చేస్తే వచ్చే సూత్రాన్ని $\phi^M$గా ఉంచుదాం. అంటే, ``$\lexists[x]$'' స్థానంలో ``$(\lexists[x \in M])$''ను, ``$\lforall[x]$'' స్థానంలో ``$(\lforall[x \in M])$''ను పెట్టాలి.
\end{defn}\noindent
$\Z$లోని ప్రతి స్వయంసిద్ధ సూత్రం $\phi$కు $\ZF \vdash
\phi^{V_{\omega+\omega}}$ అని చూపవచ్చు. కానీ \olref[spine][rank]{ordsetrankalpha} ప్రకారం $\omega+\omega$, $V_{\omega+\omega}$లో \emph{లేదు}. కనుక $\omega+\omega$ ఉనికిలో లేకపోవడం $\Z$తో సుసంగతమే.

అందుకే \olref[ordinals][replacement]{sec}లో, ప్రతిస్థాపన లేకుండా \olref[ordinals][ordtype]{thmOrdinalRepresentation}ను నిరూపించలేమని చెప్పాం. ఎందుకంటే, సహజంగా $\omega+\omega$ క్రమరకం \emph{ఉండవలసిన} స్పష్టమైన సుక్రమాన్ని $\Z$లో సులభంగా నిర్వచించవచ్చు. నిజానికి, సహజ సంఖ్యలపై ఈ కింది క్రమాన్ని ఇచ్చినప్పుడు \olref[ordinals][idea]{sec}లో దానికి అనౌపచారిక ఉదాహరణ ఇచ్చాం:
\begin{align*}
	n \lessdot m \text{ అప్పుడూ అప్పుడే }&\text{ఒకవేళ }n < m\text{ మరియు }m-n\text{ సరి సంఖ్య అయితే,}\\
	& \text{లేదా $n$ సరి సంఖ్య, $m$ బేసి సంఖ్య అయితే.}
\end{align*}
కానీ $\omega+\omega$ ఉనికిలో లేకపోతే, ఈ సుక్రమం ఏ క్రమసంఖ్యకూ సమరూపం కాదు. కనుక $\Z$, \olref[ordinals][ordtype]{thmOrdinalRepresentation}ను \emph{నిరూపించదు}.

మరో వైపు చూస్తే: ప్రతిస్థాపనతో $\omega+\omega$ ఉనికిని, దాని వల్ల $V_{\omega+\omega}$ ఉనికిని నిరూపించవచ్చు. అంతటితో కాదు. మనం నిర్వచించగల \emph{ఏ} సుక్రమానికైనా, దానికి సమరూపమైన ఏదో $\alpha$ ఉందని \olref[ordinals][ordtype]{thmOrdinalRepresentation} చెబుతుంది; కాబట్టి $V_\alpha$ కూడా ఉంటుంది. ఇలా ప్రతిస్థాపన సమితుల స్థాయిక్రమం \emph{చాలా ఎత్తుగా} ఉండాలని నిర్ధారిస్తుంది.

తదుపరి కొన్ని విభాగాల్లో, మళ్లీ \olref[card-arithmetic][fix]{sec}లో, ప్రతిస్థాపన స్థాయిక్రమాన్ని ఎంత \emph{ఎత్తుకు} తీసుకెళ్తుందో మరింత స్పష్టమవుతుంది. ప్రస్తుతానికి సరళ విషయం: ప్రతిస్థాపనకు నిజంగానే \emph{సమర్థన అవసరం}!

\end{document}
